\documentclass[17pt]{flash}
\startdatakeys
\cycle{3}
\competencesDuSocle{CA}
\connaissancesRequises{}
\competencesTravaillees{}
\motsCles{division euclidienne, quotient, reste, interprétation}
\enddatakeys
\begin{document}

\flashvisibletimertrue
\FlashSetItemSep{30pt}
\FlashSetVtop{-2em}
\FlashSetTitle{Divisions avec reste}

\FlashShowConsignes{%
    \item Répondre à chaque question en interprétant la division.
}

\flashlist{
    \flashitem{45}{On forme des équipes de 5 avec 47 élèves. Combien d'équipes complètes peut-on faire et combien d'élèves restent ?}
    \flashitem{45}{Un ruban de 83 cm est découpé en morceaux de 9 cm. Quelle longueur de ruban reste-t-il ?}
    \flashitem{45}{On range 58 livres dans des boîtes de 8 livres. Combien de boîtes faut-il prévoir pour tout ranger ?}
    \flashitem{45}{On place 93 chaises en rangées de 12. Combien de rangées complètes peut-on former ?}
}

\end{document}
